\documentclass[11pt]{article}
\usepackage[margin=1in]{geometry}
\usepackage{hyperref}
\usepackage{enumitem}
\title{Macroecological Patterns of Thermal Tolerance: A Quantitative Synthesis\\
\large A Literature Review Generated with AI Assistance}
\author{LitReview AI \and A. L. Robles Fernández\\[4pt]
\small LitReview \textbar{} ai.litreview.org\\
\small Identifier: 2608.00004}
\date{2026-08-25}

\begin{document}
\maketitle

\begin{abstract}
Thermal tolerance traits underpin predictions of species vulnerability to climate warming, yet macroecological generalizations remain contested. This review quantitatively synthesizes reported patterns in thermal limits across latitude, body size, and clades. We evaluate evidence for climate variability hypotheses explaining latitudinal gradients in thermal breadth, assess phylogenetic signal in tolerance traits, and examine mismatches between critical limits measured in the laboratory and temperatures experienced in microhabitats. We conclude that trait-based forecasts improve substantially when acclimation capacity, behavioral thermoregulation, and ontogenetic variation are modeled explicitly, and we outline minimum reporting standards for comparative thermal datasets.

\noindent\textbf{Keywords:} thermal tolerance, macroecology, climate change, latitudinal gradients.
\end{abstract}

\section*{Placeholder Notice}
This document is a placeholder entry in the LitReview repository. The full
review text will be generated with AI assistance from primary literature and
released after curation. This page establishes the identifier, metadata, and
distribution formats (LaTeX, PDF, HTML) for the entry.

\section*{Scope}
The forthcoming review will synthesize recent literature on the topic above,
covering its principal methods, findings, open challenges, and directions for
future work within the field of Ecology \& Evolution.

\section*{Data and Code Availability}
No new data or code are associated with this placeholder.

\bibliographystyle{plain}
\begin{thebibliography}{9}
\bibitem{litreview2026} LitReview Consortium (2026). \emph{LitReview: an arXiv-style repository of AI-generated literature reviews for the basic sciences.} ai.litreview.org.
\end{thebibliography}

\end{document}
